\documentclass[12pt,a4paper,oneside]{article}
\usepackage[brazil]{babel}
\usepackage[utf8]{inputenc}
\usepackage{olymp}

\setcounter{problem}{4}

\begin{document}

\begin{problem}{Equação de segundo grau (reloaded!)}{equacao2}{2}
 
Faça um programa que calcule a(s) raíz(es) de uma equação de segundo grau, se houver.
%\Notes
%
%Observações, se tiver.

\InputFile

A entrada contém três valores reais, $A$, $B$ e $C$, representando os coeficientes da equação de segundo grau. Restrições: $-10000 \leq A,B,C \leq 10000$.

%\Notes
%Neste problema, é garantido que sempre haverá duas raízes reais.


\OutputFile

Seu programa deve gerar apenas uma linha de saída, e há três situações possíveis:

\begin{enumerate}
    \item existe apenas uma raiz real: imprimir seu valor
    \item não existe raiz real: imprimir a mensagem ``\texttt{Nao ha raiz real}''
    \item existem duas raízes distintas: imprimir as raízes $x'$ e $x''$, separadas por espaço em branco, sendo $x' = \frac{-b-\sqrt{b^2 -4ac}}{2a}$ e $x'' = \frac{-b+\sqrt{b^2 -4ac}}{2a}$
\end{enumerate}

Valores devem ser impressos com 2 casas decimais e a mensagem não deve conter acentos.

%O primeiro valor é a raiz  $\frac{-b+\sqrt{\Delta}}{2a}$ e o segundo a raiz $\frac{-b-\sqrt{\Delta}}{2a}$, sendo $\Delta = b^2-4ac$.

\Example

\begin{example}
\exmp{
1 4 -4
}{
-4.83 0.83
}%
\end{example}

\begin{example}
\exmp{
1 -4 4
}{
2.00
}%
\end{example}

\begin{example}
\exmp{
1 2 3
}{
Nao ha raiz real
}%
\end{example}


%Comentários sobre o exemplo, se necessário.

\end{problem}

\end{document}
